\section{Cadre méthodologique}
\label{sec:methodology}

\subsection{Moteur de backtest}

L'ensemble des calculs repose sur un moteur de backtest vectorisé de qualité institutionnelle, capable de traiter les grilles de recherche à l'échelle de la minute (environ trois millions de barres par paire). Les portefeuilles sont calculés avec partage de cash entre paires, de sorte que le levier reflète l'exposition agrégée et non la somme naïve des positions individuelles.

L'exécution simulée est délibérément conservative~: slippage de 15~points de base sur les stratégies intraday, 10~points de base sur les stratégies journalières, fees de 5~points de base par trade. Ces hypothèses sont plus sévères que ce qu'un broker retail moderne applique en production.

\subsection{Sources de données}

\begin{center}
\renewcommand{\arraystretch}{1.2}
\begin{tabular}{lll}
\toprule
\textbf{Source} & \textbf{Fréquence} & \textbf{Usage} \\
\midrule
EUR-USD, GBP-USD, USD-JPY, USD-CAD & minute & Moteur~1 (close et VWAP intraday) \\
Idem, resamplées & journalière & Moteurs~2 et~3 (signaux EMA et RSI) \\
Spread 10Y--2Y des bons du Trésor US & journalière & Filtre macro du Moteur~1 \\
Taux de chômage US & mensuelle & Filtre macro du Moteur~1 \\
\bottomrule
\end{tabular}
\end{center}

Les séries FX minute sont strictement intersectées entre les paires (pas de \emph{forward-fill}), ce qui garantit l'absence de biais silencieux dû à un \emph{gap} sur une paire unique.

\subsection{Découpage in-sample et hors échantillon}

\begin{itemize}
\item \textbf{In-sample}~: du 1\textsuperscript{er}~avril 2019 au 1\textsuperscript{er}~avril 2025. Période utilisée pour le design de la stratégie, la sélection des paires, l'allocation et le tuning des hyperparamètres.
\item \textbf{Hors échantillon (OOS)}~: du 1\textsuperscript{er}~avril 2025 au 1\textsuperscript{er}~avril 2026. Période réservée à la validation, jamais utilisée pendant le design.
\end{itemize}

\begin{warningbox}
\textbf{Point de transparence sur le caractère de l'OOS~:} la fenêtre hors échantillon est un découpage \emph{ex-post}, effectué lorsque les données 2025-04 et postérieures sont devenues disponibles. Il ne s'agit donc pas d'un véritable \emph{holdout} pré-engagé avant le design. Malgré cette limite, le résultat reste significatif parce que le Sharpe OOS est strictement supérieur à l'in-sample ($1.44$ contre $0.94$)~--- un surajustement authentique produirait systématiquement une dégradation hors période de conception, pas une amélioration.
\end{warningbox}

\subsection{Contrôle anti look-ahead}

Une discipline anti look-ahead est appliquée systématiquement à trois niveaux~:

\begin{enumerate}
\item \textbf{Au niveau de chaque signal individuel}~: chaque indicateur (EMA, RSI, Bollinger Bands) est lagué d'un bar par rapport à sa consommation, de sorte que le signal à l'instant $t$ n'utilise que l'information strictement disponible avant $t$.
\item \textbf{Au niveau du filtre macro}~: le taux de chômage, publié le premier vendredi du mois suivant, est propagé dans le temps en respectant sa date de publication réelle. Le spread de courbe des taux est quotidien et ne nécessite pas de lag supplémentaire.
\item \textbf{Au niveau du ciblage de volatilité}~: le levier appliqué à l'instant $t$ utilise uniquement les volatilités réalisées strictement antérieures à $t$.
\end{enumerate}

Cette discipline est validée par des tests automatisés qui vérifient, pour chaque signal, que décaler les returns futurs d'un jour ne modifie pas les trades du backtest~--- un test standard de non-contamination. La reproductibilité numérique du pipeline est garantie par un test de référence qui détecte immédiatement toute dérive des résultats.
